% modules/1_introduccion.tex
% ================================================================
% CAPÍTULO 1 - INTRODUCCIÓN Y CONTEXTO DEL MANTENIMIENTO
% Referencia SWEBOK V4.0a, Capítulo 7: Software Maintenance
% ================================================================

\chapter{Introducción}

\section{Propósito del Documento}

El presente documento establece el Plan de Mantenimiento de Software para el sistema
educativo \textbf{Yoliztli}, una plataforma web interactiva orientada a la enseñanza de
lenguas originarias --- náhuatl y otomí --- para niños de educación primaria en el estado
de Hidalgo.

El plan se rige bajo los lineamientos del estándar internacional \textbf{ISO/IEC/IEEE 14764}
y la guía de conocimiento \textbf{SWEBOK V4.0a (Capítulo 7: Software Maintenance)}, los
cuales definen las actividades, categorías y procesos requeridos para mantener un producto
de software en operación continua y satisfactoria.

\section{Alcance}

Este documento aplica a la versión del sistema Yoliztli gestionada como \textit{fork} del
repositorio original. El alcance del mantenimiento comprende:

\begin{itemize}
    \item La corrección de defectos reportados en el ciclo de pruebas previo.
    \item La transición de infraestructura tecnológica (motor de base de datos).
    \item La incorporación de nuevas funcionalidades requeridas por el SRS original no
          implementadas en la versión anterior.
    \item La mejora interna de la calidad y mantenibilidad del código.
\end{itemize}

\section{Contexto del Proyecto}

\subsection{Descripción General del Sistema}

Yoliztli es una aplicación web construida sobre el stack tecnológico siguiente:

\begin{itemize}
    \item \textbf{Backend:} Laravel 12 (PHP 8.2+), siguiendo el patrón de arquitectura MVC.
    \item \textbf{Frontend:} React 18 integrado mediante Inertia.js como adaptador SPA.
    \item \textbf{Base de Datos:} SQLite (desarrollo local) y PostgreSQL (producción).
    \item \textbf{Compilación:} Vite 5 con plugin de Laravel y React.
    \item \textbf{Estilización:} Tailwind CSS 3.
\end{itemize}

\subsection{Antecedentes de la Versión Anterior}

La versión original del sistema planteaba el uso de \textbf{Firebase} como plataforma
backend (\textit{Serverless}), con autenticación, base de datos en tiempo real y
almacenamiento en la nube. El equipo de desarrollo del \textit{fork} actual realizó una
migración arquitectónica completa hacia una solución autohospedada con Laravel y SQLite,
lo que otorga mayor independencia del entorno, control del flujo de datos y
compatibilidad con infraestructura institucional propia.

\subsection{Estado de las Pruebas al Inicio del Mantenimiento}

El reporte de calidad (\textit{5.Pruebas Yoliztli.pdf}) ejecutado el 9 de diciembre de 2025
registra los siguientes resultados sobre 8 casos de prueba totales:

\begin{table}[H]
\centering
\begin{tabular}{clc}
\toprule
\textbf{ID} & \textbf{Caso de Prueba} & \textbf{Resultado} \\
\midrule
CP001 & Registro de usuario (nombre, correo, contraseña)  & \textcolor{process}{Aprobado} \\
CP002 & Inicio de sesión                                  & \textcolor{process}{Aprobado} \\
CP003 & Acceso a recursos interactivos                    & \textcolor{process}{Aprobado} \\
CP004 & Selección de idioma de interfaz                   & \textcolor{process}{Aprobado} \\
CP005 & Panel de reportes de docentes                     & \textcolor{process}{Aprobado} \\
CP006 & Descarga de materiales educativos                 & \textcolor{process}{Aprobado} \\
CP007 & Verificar filtro por categoría y edad             & \textcolor{red}{Fallido}     \\
CP008 & Verificar el avance/progreso del usuario          & \textcolor{red}{Fallido}     \\
\bottomrule
\end{tabular}
\caption{Resultados del ciclo de pruebas de aceptación --- Diciembre 2025}
\label{tab:resultados-pruebas}
\end{table}

Los casos CP007 y CP008 representan la base de las actividades de
\textbf{mantenimiento correctivo y aditivo} planificadas en este documento.
